% Auto-generated by src/make_tables.py -- do not edit by hand.
\begin{table}[t]
\centering
\caption{Multilevel Picard, measured rather than cited. \emph{Upper}: the method realises the Picard iteration, so it needs that iteration to be a contraction, with Lipschitz constant $|\nabla V|/\eta$. Sweeping $\eta$ from weak control authority to strong locates the boundary between $|\nabla V|/\eta = 0.56$ (converged) and $0.91$ (diverged) instead of leaving it as a caveat. \emph{Lower}: the same HJB family at $d = 1$ to $20$. At the paper's own calibration the scheme diverges at every dimension --- in a tenth of a second, so this is the hypothesis failing rather than a budget running out --- while at $\eta = 1$ it is the cheapest method in this paper, $0.26$\% at $d=20$ in $2.3$\,s on a CPU, where no grid can run and where the PINN reaches $R^2 = 0.757$. The implementation is checked where the hypothesis holds: at $\eta=1$ the relative error falls $6.7e-03 \to 1.5e-03$ from level 2 to level 4.}
\label{tab:mlp}
\footnotesize
\setlength{\tabcolsep}{4pt}
\begin{tabular}{rrrr}
\toprule
\multicolumn{4}{l}{\emph{contraction boundary at $d=1$}} \\
$\eta$ & $|\nabla V|/\eta$ & level-3 rel.\ err & verdict \\
\midrule
$10$ & $0.011$ & $7.02e-03$ & converged \\
$1$ & $0.100$ & $6.37e-03$ & converged \\
$0.3$ & $0.275$ & $1.06e-02$ & converged \\
$0.1$ & $0.557$ & $2.12e-02$ & converged \\
$0.03$ & $0.907$ & $4.38e+00$ & diverged \\
$0.01$ & $1.249$ & $8.21e+02$ & diverged \\
\midrule
\multicolumn{4}{l}{\emph{across dimension, $M=2$, 3 levels}} \\
$d$ & rel.\ err at $\eta=0.01$ & rel.\ err at $\eta=1$ & time at $\eta=1$ (s) \\
\midrule
$1$ & $82115$\% & $0.637$\% & $0.00$ \\
$4$ & $582833$\% & $0.439$\% & $0.05$ \\
$8$ & $551822$\% & $0.425$\% & $0.20$ \\
$12$ & $433216$\% & $0.173$\% & $0.58$ \\
$20$ & $201008$\% & $0.258$\% & $2.25$ \\
\bottomrule
\end{tabular}
\end{table}
